\documentclass[12pt]{article}
\usepackage[utf8]{inputenc}
\usepackage[T1]{fontenc}
\usepackage[margin=1in]{geometry}
\usepackage{setspace}
\usepackage{xcolor}
\usepackage[colorlinks=true, urlcolor=blue!60!black, linkcolor=blue!60!black]{hyperref}
\setlength{\parindent}{2em}
\setlength{\parskip}{1.25em}
\renewcommand{\baselinestretch}{1.0}
\begin{document}
\begin{center}
{\LARGE \textit{Michael H. v. Gerald D.}, 491 U.S. 110 (1989)\par}
\vspace{0.8em}
{\large Case Brief --- Four-Page Edition\par}
\end{center}
\vspace{1em}

\section{Memory Jogger}

A fractured Court upheld California's conclusive presumption that a child born to a married woman cohabiting with her husband is a child of the marriage, rejecting a biological father's due process claim and sparking a landmark debate over whether substantive due process rights must be defined at the most specific level of historical tradition (\href{https://supreme.justia.com/cases/federal/us/491/110/}{491 U.S. 110}).

\section{Detailed Case Facts}

Carole D. married Gerald D. in 1976 and began an affair with a neighbor, Michael H., in 1978. Victoria D. was born in May 1981; Gerald was listed as father and held her out as his daughter. Blood tests showed a 98.07\% probability that Michael was the biological father. Over the next three years Carole and Victoria lived intermittently with Michael, with another man, and with Gerald. From August 1983 to May 1984 they lived with Michael, who supported Victoria and whom she called ``Daddy''; Carole signed but never filed a stipulation acknowledging his paternity. In June 1984 she reconciled permanently with Gerald in New York.

California Evidence Code \S~621 provided that ``the issue of a wife cohabiting with her husband, who is not impotent or sterile, is conclusively presumed to be a child of the marriage,'' rebuttable only on a timely motion by the husband or, in limited cases, the mother. Civil Code \S~4601 allowed discretionary visitation for ``any other person having an interest in the welfare of the child.''

\section{Procedural History}

Michael filed a filiation action in 1982; Victoria's guardian ad litem claimed a right to relationships with both fathers. After granting Michael limited visitation, the Superior Court granted Gerald summary judgment under \S~621 and rejected the constitutional challenges. The Court of Appeal affirmed, the California Supreme Court denied review, and the U.S. Supreme Court affirmed on June 15, 1989.

\section{Judicial Votes}

Affirmed, \textbf{5--4}, with no majority opinion.
\begin{itemize}
\item \textbf{Plurality:} Justice Scalia, joined by Chief Justice Rehnquist and, except Footnote 6, by Justices O'Connor and Kennedy.
\item \textbf{Concurring in part:} Justice O'Connor, joined by Justice Kennedy, rejecting Footnote 6.
\item \textbf{Concurring in the judgment:} Justice Stevens.
\item \textbf{Dissents:} Justice Brennan, joined by Justices Marshall and Blackmun; Justice White, joined by Justice Brennan.
\end{itemize}

\section{Holding}

A State does not violate due process by applying a conclusive presumption of marital paternity to bar a biological father from establishing paternity of a child conceived and born into an extant, cohabiting marriage whose spouses wish to raise the child. The child has no due process right to dual fatherhood, and limiting rebuttal to the spouses survives rational-basis review. Footnote 6's ``most specific level'' method drew only two votes and was rejected by seven Justices; Footnote 7's reservation for marital parents who do not wish to raise the child is also dicta.

\section{Analysis of Opinions}

\textbf{Scalia (plurality).} Section 621 is a substantive rule of law, not a procedural presumption: California has deemed an outsider's paternity irrelevant to a child born into an existing marriage. A protected liberty interest must be fundamental and ``traditionally protected by our society.'' The unwed-father cases (\textit{Stanley}, \textit{Quilloin}, \textit{Caban}, \textit{Lehr}) rest on respect for the ``unitary family,'' and from Bracton and Blackstone onward tradition protected the marital family against claims by an ``adulterous natural father.''

\textbf{Footnote 6 (Scalia and Rehnquist).} Courts must consult the most specific tradition protecting or denying the asserted right.

\textbf{O'Connor (concurring in part).} A single ``most specific level'' test ``may be somewhat inconsistent with our past decisions,'' including \textit{Griswold}, \textit{Eisenstadt}, and \textit{Loving}.

\textbf{Stevens (concurring in the judgment).} Assuming Michael had a liberty interest, \S~4601's visitation avenue gave him a fair opportunity to be heard.

\textbf{Brennan (dissenting).} Five Justices refused to foreclose a natural father's interest, and only two endorsed Footnote 6, under which \textit{Griswold}, \textit{Eisenstadt}, and \textit{Stanley} would have come out differently. Biology plus a developed relationship creates a liberty interest; Stevens's reading of \S~4601 was ``mere wishful thinking.''

\textbf{White (dissenting).} Michael satisfied \textit{Lehr}: he lived with Victoria, supported her, and held her out as his child. Blood tests removed old evidentiary doubts, yet \S~621 barred him without a hearing.

\section{Examples: Future Applications}

\begin{enumerate}
\item \textbf{Same-Side (Upheld): Intact Marriage.} A couple living together wishes to raise their child; the wife's former affair partner seeks DNA testing and custody three years later. The marital presumption controls.
\item \textbf{Same-Side (Upheld): Three Parents.} A child claims a due process right to legal parentage with her mother, the mother's husband, and her biological father. Tradition does not require recognizing multiple fathers.
\item \textbf{Opposite-Side: Unmarried Mother.} An unwed father who raised his children for six years loses them to foster care after the mother dies, under a presumption of unfitness. With no marital family, \textit{Stanley} requires a hearing.
\item \textbf{Opposite-Side: Dissolved Marriage.} The wife has left her husband, who disclaims paternity, and lives with the biological father. Footnote 7 limited the holding to cohabiting spouses who wish to raise the child.
\item \textbf{Fence-Sitter: Known Donor.} A married same-sex couple invites a known donor into the child's life, then cuts him off. \textit{Michael H.} favors the marital unit, but a co-parenting invitation and the absence of a \S~4601-style avenue could distinguish it.
\end{enumerate}

\section{Critique}

\textbf{Originalist and textualist views.} Defenders argue that because the Due Process Clause guarantees ``process,'' any unenumerated rights must be disciplined by concrete historical practice, and consulting the most specific tradition keeps judges from smuggling in personal preferences under labels like ``family.''

\textbf{Progressive and doctrinal views.} Laurence Tribe and Michael Dorf showed that choosing a level of generality is never value-neutral; even the ``most specific'' tradition requires choices about which facts matter. Critics add that tradition fails to constrain judges, and family-law scholars fault the plurality for elevating marital form over the lived attachments of child and parent.

\section{Key Quotations}

\begin{quote}
``As we view them, they rest not upon such isolated factors but upon the historic respect --- indeed, sanctity would not be too strong a term --- traditionally accorded to the relationships that develop within the unitary family.''
\end{quote}
\noindent --- Justice Scalia, plurality opinion, 491 U.S. at 123.

\begin{quote}
``We refer to the most specific level at which a relevant tradition protecting, or denying protection to, the asserted right can be identified. ... Although assuredly having the virtue (if it be that) of leaving judges free to decide as they think best when the unanticipated occurs, a rule of law that binds neither by text nor by any particular, identifiable tradition is no rule of law at all.''
\end{quote}
\noindent --- Justice Scalia, plurality opinion (Footnote 6, joined only by Chief Justice Rehnquist), 491 U.S. at 127--128 n.6.

\begin{quote}
``The document that the plurality construes today is unfamiliar to me. It is not the living charter that I have taken to be our Constitution; it is instead a stagnant, archaic, hidebound document steeped in the prejudices and superstitions of a time long past.''
\end{quote}
\noindent --- Justice Brennan, dissenting, 491 U.S. at 141.

\end{document}
